% META: {"id":"1SPE-EXPO-CO-045","chapitre":"1SPE-EXPONENTIELLE","type_objet":"corrige","exercice_id":"1SPE-EXPO-EX-045","status":"generated"}
% BEGIN-VERIFY
% from sympy import *
% n = symbols('n', integer=True, nonnegative=True)
% u = 120*exp(Rational(1, 20)*n)
% assert simplify(u.subs(n, n + 1)/u - exp(Rational(1, 20))) == 0
% assert round(float(u.subs(n, 5)), 1) == 154.1
% END-VERIFY
\begin{corrige}{1SPE-EXPO-EX-045}
\begin{enumerate}
\item $u_n=Q(n)=120\,\mathrm{e}^{0{,}05n}$.
\item
\[\frac{u_{n+1}}{u_n}=\frac{120\,\mathrm{e}^{0{,}05(n+1)}}{120\,\mathrm{e}^{0{,}05n}}=\mathrm{e}^{0{,}05}.\]
Le quotient étant constant, $(u_n)$ est géométrique de raison $\mathrm{e}^{0{,}05}$.
\item $Q(5)=120\,\mathrm{e}^{0{,}25}\approx154{,}1$ et, par définition, $u_5=Q(5)$. Le modèle discret coïncide avec le modèle continu aux dates entières.
\item Le coefficient $120$ est la valeur initiale. D'une date entière à la suivante, la quantité est multipliée par $\mathrm{e}^{0{,}05}\approx1{,}051$.
\end{enumerate}
\end{corrige}
